\documentclass[spanish,letterpaper,oneside]{article}
\def\documenttitle {Delitos contra la vida y la integridad corporal}
\def\documentsubtitle {Vida, integridad y subsunción penal}
\def\documentsubject {Actividad 3 - Derecho penal especial mexicano}
\def\documentauthor {Martín Jonathan de la Cruz}
\def\coursename {Derecho penal especial mexicano}
\def\coursecode {020}
\def\universityname {Universidad Abierta y a Distancia de México}
\def\universityfaculty {Licenciatura en Derecho}
\def\universitydepartment {Derecho penal especial mexicano}
\def\universitydepartmentimage {departamentos/UnADM}
\def\universitydepartmentimagecfg {height=1.57cm}
\def\universitylocation {Roma Norte, Ciudad de México}
\def\authortable {
	\begin{tabular}{ll}
		\textbf{Alumno:} & Martín Jonathan de la Cruz \\
		Matrícula: & ES2611202040 \\
		Figura docente: & Verónica Ortiz López \\
		Semestre/Bloque: & 2 / 2 \\
		Semana: & 3 \\
		\multicolumn{2}{l}{Fecha: 12 de octubre de 2026} \\
		\multicolumn{2}{l}{\universitylocation}
	\end{tabular}
}
\input{template}
\def\indexstyle{}
\usepackage{helvet}
\renewcommand{\familydefault}{\sfdefault}
\usepackage[most]{tcolorbox}
\usepackage{pdflscape}
\usepackage{tikz}
\usetikzlibrary{babel}
\setcitestyle{authoryear,open={(},close={)}}
\newcommand{\insertcoverwatermark}{%
	\AddToShipoutPictureBG*{%
		\begin{tikzpicture}[remember picture,overlay]
			\node[opacity=0.16,inner sep=0pt,yshift=-0.35cm]
				at (current page.center) {\includegraphics[width=0.72\paperwidth]{img/departamentos/UnADM.pdf}};
		\end{tikzpicture}%
	}%
}
\definecolor{casoverde}{HTML}{174A3A}
\definecolor{casopapel}{HTML}{F6F7F2}
\newtcolorbox{casoestudio}[1]{
	enhanced,breakable,colback=casopapel,colframe=casoverde,
	boxrule=0pt,leftrule=2.5pt,arc=1mm,
	left=3mm,right=3mm,top=2mm,bottom=1.5mm,fontupper=\small,
	title={#1},colbacktitle=casoverde,coltitle=white,
	fonttitle=\bfseries,before skip=8pt,after skip=8pt
}
\newcommand{\fichacaso}[3]{%
	\begin{tcolorbox}[colback=casoverde!5,colframe=casoverde!30,
		boxrule=0.4pt,arc=0.6mm,left=2mm,right=2mm,top=1mm,bottom=1mm]
	\textbf{Hechos:} #1\par
	\textbf{Problema jurídico:} #2\par
	\textbf{Normas y criterio:} #3
	\end{tcolorbox}%
}
\setlength{\parindent}{0pt}
\setlength{\parskip}{3pt}
\renewcommand{\arraystretch}{1.13}
\sloppy
\begin{document}
\insertcoverwatermark
\templatePortrait
\templatePagecfg
\clearpage
\begin{abstractd}
El resultado visible de una agresión no identifica por sí solo el delito: la muerte exige nexo causal e imputación; la fractura, clasificación médica; y ciertas figuras requieren además parentesco conocido, razones de género o una forma concreta de intervención. El encuadramiento debe confrontar hechos disponibles, elementos legales, hipótesis alternativas y pruebas capaces de confirmarlas o descartarlas. El Código Penal Federal y el texto consolidado de Ciudad de México presentan diferencias que no deciden por sí mismas la competencia. Los amparos en revisión 554/2013 y 267/2023 orientan, respectivamente, la investigación de muertes violentas de mujeres y el examen constitucional de la criminalización federal del aborto. Un juicio jurídico responsable protege a víctimas y derechos reproductivos sin convertir indicios en culpabilidad probada.

\textbf{Palabras clave:} homicidio; lesiones; feminicidio; violencia de pareja; participación en el suicidio; aborto; subsunción penal.
\end{abstractd}
\templateIndex
\templateFinalcfg
% Técnica #66: cada caso conserva contexto, problema, objetivo, desarrollo y reflexión.
% Técnica #1: el cuadro contrasta cinco casos por categorías comunes.
\section{Introducción}
Cinco situaciones sobre lesiones, muerte, violencia de pareja, suicidio y aborto plantean un mismo problema: cómo encuadrar hechos incompletos sin convertir hipótesis en condenas. El lugar de los hechos no consta, salvo el servicio médico del caso 5 en Ciudad de México: los artículos del Código Penal para el Distrito Federal (CPDF) permiten formular hipótesis locales; los del Código Penal Federal (CPF) sirven de contraste, sin asignar de oficio competencia federal. Distinguir una \emph{hipótesis típica} de una responsabilidad demostrada requiere verificar conducta, sujetos, resultado, nexo, elemento subjetivo y, en su caso, circunstancias calificativas \citep[cap.~3]{amuchateguiPenal2012}. La Constitución (arts.~1 y 4) orienta la interpretación protectora de derechos sin sustituir el tipo penal \citep{cpeumAulaS3}.

\section{Encuadramiento de delitos contra la vida y la integridad corporal}
La protección de la vida y la integridad comprende figuras con bienes y requisitos distintos \citep[parte especial, sección 1, caps.~13--20]{amuchateguiPenal2012}. Amuchategui Requena separa sujeto activo, persona afectada, conducta, bien protegido y las calidades o conocimientos que un tipo puede exigir adicionalmente \citep[cap.~3, pp.~38--40]{amuchateguiPenal2012}. La doctrina de 2012 explica esas categorías; las formulaciones vigentes requieren el contraste de códigos consolidados \citep{cpfOficial2026,cpdfS3Consolidado}. La diferencia entre resultado, intención y circunstancia especial evita convertir una agresión en una calificativa por intuición.

\subsection*{Homicidio y lesiones: resultado, sujetos y nexo causal}
\phantomsection\addcontentsline{toc}{subsection}{Homicidio y lesiones: resultado, sujetos y nexo causal}
\textbf{Homicidio y causalidad.} El CPDF 123 y el CPF 302 exigen la privación de la vida de otra persona; el sujeto activo del supuesto básico no necesita parentesco, mientras el pasivo es quien fallece. La muerte posterior a una agresión no equivale automáticamente a homicidio: CPDF 124 y CPF 303--305 obligan a examinar si la lesión produjo el desenlace o intervino una causa ajena. El dictamen forense, la cronología y el elemento subjetivo son variables distintas. En CPDF 138, ventaja, alevosía u otras calificativas exigen condiciones propias; la riña del 129 no surge de un mero intercambio verbal \citep{cpdfS3Consolidado,cpfOficial2026}.

\textbf{Lesiones y su gravedad.} El CPF 288 comprende fracturas y alteraciones de salud causadas externamente; CPDF 130 clasifica la lesión por tiempo de sanación, secuela, pérdida funcional o peligro para la vida. El activo causa el daño y el pasivo lo resiente; se protege la salud e integridad corporal, no un plazo supuesto por intuición. La fractura requiere diagnóstico y pronóstico médico para elegir fracción; CPDF 131 y 138 contemplan consecuencias adicionales solo si se prueban sus presupuestos. Lesión y homicidio comparten la agresión posible, pero difieren en resultado e imputación \citep{cpfOficial2026,cpdfS3Consolidado}.

\clearpage
\phantomsection\addcontentsline{toc}{subsection}{Caso 1. Lesiones y homicidio}
\begin{casoestudio}{Caso 1. Lesiones de Pedro y muerte de otra persona}
\fichacaso{Carlos golpea a Pedro y le causa una fractura. En un hecho distinto Roberto agrede a otra persona que después fallece; el supuesto no identifica lugar ni mecanismo de muerte.}{Distinguir lesión de homicidio, sujetos y bienes protegidos, y precisar qué datos permitirían apreciar agravantes o atenuantes.}{CPDF, arts.~123--124, 129--130 y 138; CPF, arts.~288, 302--305 y 315.}
\textit{La fractura y la muerte exigen encuadramientos distintos.} Carlos golpea voluntariamente a Pedro y le causa una fractura, sin que se indique fallecimiento: la hipótesis es lesiones (CPDF, art.~130; CPF, art.~288). La fractura no establece por sí sola días de sanación, secuelas ni peligro vital; se necesita dictamen. Roberto agrede a otra persona que muere después: si el fallecimiento se debe a la lesión y se prueban los requisitos de imputación, procede investigar homicidio (CPDF, arts.~123--124; CPF, arts.~302--305). La tesis alternativa es que la muerte se explique por causa independiente o no resulte atribuible en la modalidad de homicidio propuesta; el mero orden temporal no prueba nexo ni dolo de matar \citep{cpdfS3Consolidado,cpfOficial2026}.

\textit{Los sujetos, resultados y bienes permiten separar los hechos.} Carlos y Roberto son posibles sujetos activos; Pedro y la persona fallecida, pasivos. La acción es agredir; los resultados descritos son, respectivamente, fractura y muerte. El primer bien es salud e integridad corporal; el segundo, vida. El dolo de golpear no equivale necesariamente a dolo de causar la muerte.

\textit{Las modificativas requieren hechos adicionales.} CPDF 138 contempla ventaja, traición, alevosía, retribución, medio empleado, saña, alteración voluntaria u odio, con requisitos específicos; CPDF 129 trata la riña. Una discusión no acredita ninguna de ellas: tampoco una contienda física recíproca. Investigar arma, mecánica, fuerza, defensas, provocación y testimonios; solicitar expediente clínico, necropsia y peritaje causal. El CPF 315 regula otras calificativas; no trasplantar agravantes entre fueros. La solución provisional es tratar el primer hecho como lesiones y abrir para el segundo la hipótesis de homicidio, condicionada al nexo y al elemento subjetivo \citep{cpdfS3Consolidado,cpfOficial2026}.
\end{casoestudio}

\subsection*{Parentesco y conocimiento del vínculo}
\phantomsection\addcontentsline{toc}{subsection}{Parentesco y conocimiento del vínculo}
CPDF 125 y CPF 323 añaden a la muerte la relación legalmente prevista y el conocimiento de ella por el activo. Padre e hijo satisfacen en abstracto la calidad de ascendiente y descendiente en línea recta, pero la filiación y su conocimiento son cuestiones de prueba; el dolo de agredir tampoco decide por sí solo el de matar. Sin conocimiento, ambos códigos remiten a reglas de homicidio simple, sin borrar el examen del nexo causal \citep{cpdfS3Consolidado,cpfOficial2026}.

El alcance de la figura especial depende del vínculo descrito en cada ley y no de una noción genérica de ``familia'': CPDF 125 incluye hijastro y pareja permanente; CPF 323 menciona convivente y compañera o compañero civil. Padre e hijo son ascendiente y descendiente en ambos, pero no deben intercambiarse listas ni penas. Para atribuir el resultado, el registro de filiación debe relacionarse con datos de conocimiento efectivo y con la reconstrucción médica y testimonial de la agresión \citep[art.~125]{cpdfS3Consolidado}\citep[art.~323]{cpfOficial2026}.

\Needspace{12\baselineskip}
\phantomsection\addcontentsline{toc}{subsection}{Caso 2. Miguel y su padre}
\begin{casoestudio}{Caso 2. Miguel y su padre}
\fichacaso{En una discusión familiar Miguel agrede intencionalmente a su padre y este muere; Miguel conoce el parentesco.}{Identificar los elementos específicos del vínculo conocido y distinguirlos del homicidio simple sin presumir dolo de matar.}{CPDF, arts.~123, 125 y 138; CPF, arts.~302 y 323.}
\textit{La muerte del ascendiente exige probar todos los elementos del tipo.} La muerte de un ascendiente consanguíneo en línea recta, con conocimiento del vínculo, lleva a analizar CPDF 125 (o CPF 323 si se justificara su aplicación). En el supuesto se afirma filiación conocida y agresión intencional; verificar muerte, nexo, filiación, conocimiento y el elemento subjetivo relativo a la privación de la vida \citep{cpdfS3Consolidado,cpfOficial2026}.

\textit{El conocimiento del parentesco tiene efecto jurídico propio.} CPDF 125 lo exige y remite a la punibilidad del homicidio simple si falta. Un registro civil y comunicaciones son pruebas pertinentes, sin confundir parentesco probado con conocimiento automáticamente demostrado \citep[art.~125]{cpdfS3Consolidado}.

\textit{La figura especial añade requisitos al homicidio simple.} CPDF 123 requiere privar de la vida, pero no relación familiar ni saber de ella. Si el ataque buscó solo lesionar y la muerte no es imputable como homicidio doloso, no basta invocar el parentesco para cerrar el análisis. La solución fundada es investigar la figura especial y, según la prueba del nexo y del dolo, contrastarla con homicidio simple; el conflicto familiar no configura por sí solo una calificativa del CPDF 138.
\end{casoestudio}

\subsection*{Violencia de pareja y razones de género}
\phantomsection\addcontentsline{toc}{subsection}{Violencia de pareja y razones de género}
\textbf{Violencia de pareja.} El CPDF 200, reformado en diciembre de 2025, incluye acción u omisión de violencia física, psicoemocional, sexual, económica, patrimonial o contra derechos reproductivos y protege también a la expareja en las relaciones allí definidas. El CPF 343 Bis enumera dominio, control y agresiones dentro o fuera del domicilio familiar. Identificar cada episodio, su fecha, vínculo y medio de prueba permite distinguir el patrón previo del resultado letal posterior, sin decidir de manera automática un concurso ni aplicar simultáneamente ambos fueros \citep{cpdfS3Consolidado,cpfOficial2026}.

\textbf{Feminicidio y contexto.} La muerte de una mujer por razones de género es el núcleo de CPDF 148 Bis y CPF 325. En el texto local, los antecedentes de amenazas o violencia del activo (fr.~III) y la relación afectiva actual o pasada (fr.~V) son vías normativas para examinar esas razones; deben acreditarse junto con autoría y resultado. El AR 554/2013 de la Primera Sala ordena investigar las muertes violentas de mujeres con perspectiva de género, incluso cuando una hipótesis inicial apunte a suicidio; el criterio obliga a indagar, no a presumir una condena \citep{cpdfS3Consolidado,cpfOficial2026,scjnAR5542013}.

\Needspace{12\baselineskip}
\phantomsection\addcontentsline{toc}{subsection}{Caso 3. Ana y su expareja}
\begin{casoestudio}{Caso 3. Ana y su expareja}
\fichacaso{Ana sufrió amenazas, control económico y agresiones físicas de su pareja durante años. Tras la ruptura continuaron las amenazas y su expareja la privó de la vida.}{Examinar violencia de pareja, posibles razones de género y la distinción entre homicidio y feminicidio.}{CPDF, arts.~123, 148 Bis (frs.~III y V) y 200; SCJN, AR 554/2013.}
\textit{La violencia de pareja comprende actos anteriores y posteriores a la ruptura.} Las amenazas, el control económico y las agresiones pueden configurar violencia familiar entre pareja (CPDF, art.~200, reformado en 2025), incluso respecto de una expareja en las relaciones allí definidas. Deben individualizarse episodios y probarse el vínculo y cada modalidad; agresiones específicas pueden además exigir examen de lesiones o amenazas, sin duplicar automáticamente sanciones por el mismo hecho \citep[art.~200]{cpdfS3Consolidado}.

\textit{Los antecedentes y el vínculo orientan la investigación del feminicidio.} La privación de la vida de Ana por su expareja exige investigar CPDF 148 Bis: antecedentes de amenazas/violencia (fr.~III) y una relación sentimental o afectiva pasada (fr.~V) son supuestos normativos relevantes, a corroborar con denuncias, mensajes, testimonios, pericias y cronología \citep[art.~148 Bis]{cpdfS3Consolidado}. En el fuero federal, CPF 325 contiene supuestos propios, no idénticos. La Primera Sala en el amparo en revisión 554/2013 exigió investigar muertes violentas de mujeres con perspectiva de género y subsanar omisiones; no declaró que la mera condición de mujer pruebe autoría ni culpabilidad \citep{scjnAR5542013}.

\textit{Las razones de género diferencian feminicidio de homicidio.} Homicidio simple CPDF 123 exige muerte atribuible; feminicidio agrega las razones de género previstas en CPDF 148 Bis. Los antecedentes y el vínculo ameritan abrir la línea de feminicidio, pero una relación pasada sin acreditación del supuesto típico, autoría y conexión probatoria con la muerte no permite condenar. Procede investigar ambas hipótesis preservando la presunción de inocencia y los derechos de Ana; decidir el concurso con la violencia previa solo tras separar conductas y bienes afectados.
\end{casoestudio}

\subsection*{Intervención ajena en la decisión suicida}
\phantomsection\addcontentsline{toc}{subsection}{Intervención ajena en la decisión suicida}
\textbf{Participación en el suicidio.} El CPDF 142 diferencia ayudar materialmente, inducir la decisión y ejecutar uno mismo la muerte; el CPF 312 también distingue auxilio o inducción de ejecución directa. Interviene como activo quien persuade o facilita y como persona afectada quien actúa contra su vida. La autonomía y capacidad de esta última, la cronología de su intención y el aporte real del tercero evitan equiparar toda conversación con una inducción causal. La edad o falta de capacidad abren los tratamientos especiales de CPDF 143 y CPF 313; no se presumen \citep{cpdfS3Consolidado,cpfOficial2026}.

La consumación modifica el examen del auxilio: CPDF 142 distingue el suicidio consumado de la tentativa en que se causaron lesiones y del intento sin lesiones. Esa diferencia no convierte a quien entregó un medio en ejecutor de la muerte; exige reconstruir quién realizó el acto final y qué relación guardó con la colaboración. Cuando la persona ya había manifestado su intención, los mensajes y su secuencia temporal ayudan a separar una decisión anterior de una influencia posterior relevante \citep[art.~142]{cpdfS3Consolidado}.

\Needspace{12\baselineskip}
\phantomsection\addcontentsline{toc}{subsection}{Caso 4. Inducción y auxilio al suicidio}
\begin{casoestudio}{Caso 4. Inducción y auxilio al suicidio}
\fichacaso{Una persona anuncia que se quitará la vida. Un conocido la persuade, le proporciona medios y ella fallece a consecuencia de sus propios actos.}{Distinguir inducción, auxilio y ejecución directa, e identificar sujetos y elementos sin presumir el efecto de la persuasión.}{CPDF, arts.~142--143; CPF, arts.~312--313.}
\textit{Persuadir y suministrar medios son intervenciones diferentes.} El conocido persuade y aporta medios; la persona fallece por actos propios. CPDF 142 distingue \emph{inducir} (influir en la decisión) de \emph{ayudar} (facilitar su ejecución); CPF 312 regula ambos supuestos \citep{cpdfS3Consolidado,cpfOficial2026}. Su intención de suicidarse previa abre una alternativa: no basta llamar ``inducción'' a cualquier conversación posterior; el auxilio material requiere igualmente prueba de una colaboración voluntaria y relevante.

\textit{La vida de la persona fallecida es el bien protegido.} El conocido sería sujeto activo y la persona fallecida, destinataria de la inducción o auxilio. Deben documentarse su decisión previa, el contenido y oportunidad de la persuasión, la naturaleza de los medios, la conciencia del tercero, la capacidad de la persona y el nexo con el resultado. Sin esos datos no cabe afirmar que la persuasión fue determinante.

\textit{La ejecución directa no se desprende del relato.} Si el conocido hubiese causado él mismo la muerte, CPDF 142 y CPF 312 prevén un supuesto diferente y una respuesta propia; aquí la persona actuó por sí misma. Confirmar edad y capacidad antes de considerar CPDF 143 o CPF 313. La solución es investigar por separado inducción y auxilio, no sumar mecánicamente modalidades por los mismos actos \citep{cpdfS3Consolidado,cpfOficial2026}.
\end{casoestudio}

\subsection*{Interrupción del embarazo, tiempo y consentimiento}
\phantomsection\addcontentsline{toc}{subsection}{Interrupción del embarazo, tiempo y consentimiento}
\textbf{Aborto, temporalidad y derechos.} El CPDF 144 define el aborto penal a partir de la interrupción después de la semana doce; el 145 considera el consentimiento y el 146 la ausencia de este, mientras el 148-I prevé excluyente por violación. Frente a ello, CPF 329 conserva otra formulación temporal. La Primera Sala, en AR 267/2023, identificó porciones inconstitucionales de los artículos federales 330--334: el contraste impide presentar la criminalización federal como regla absoluta y tampoco convierte el amparo en una reforma textual del código local. Gestación, consentimiento y origen del embarazo requieren comprobación respetuosa de derechos \citep{cpdfS3Consolidado,cpfOficial2026,scjnAR2672023}.

Las situaciones previstas en CPDF 148 no se agotan en la violación: también contempla peligro grave para la salud, determinadas alteraciones genéticas y una conducta culposa de la embarazada, con requisitos distintos. En los supuestos que requieren valoración clínica, la información objetiva y oportuna permite una decisión libre e informada. Distinguir esas hipótesis protege la autonomía sin ocultar la importancia jurídica de las semanas de gestación ni confundir la falta de consentimiento, tratada en CPDF 146, con una solicitud voluntaria \citep[arts.~146--148]{cpdfS3Consolidado}.

\Needspace{12\baselineskip}
\phantomsection\addcontentsline{toc}{subsection}{Caso 5. Dos solicitudes en Ciudad de México}
\begin{casoestudio}{Caso 5. Dos solicitudes en Ciudad de México}
\fichacaso{Una mujer solicita interrumpir voluntariamente su embarazo de diez semanas en una institución médica de CDMX. Otra lo solicita a las dieciséis semanas porque el embarazo resultó de una violación.}{Comparar plazo, consentimiento, causa de la solicitud y disposiciones aplicables a cada supuesto.}{CPDF, arts.~144--148; CPEUM, arts.~1 y 4; SCJN, AR 267/2023.}
\textit{El plazo distingue los dos supuestos, pero no decide por sí solo su tratamiento.} A diez semanas, la interrupción consentida no cae en la definición penal temporal de CPDF 144 (después de la décima segunda semana); deben verificarse gestación y consentimiento informado. A dieciséis semanas, aunque el umbral del 144 se supera, el embarazo producto de violación encuadra, de acreditarse la circunstancia, en la excluyente de responsabilidad del CPDF 148, fr.~I; no presumir sanción por el plazo aislado \citep{cpdfS3Consolidado}.

\textit{Gestación, consentimiento y origen del embarazo son datos indispensables.} Importan las constancias clínicas de edad gestacional y consentimiento libre, el origen del embarazo y la información objetiva, suficiente y oportuna conforme al último párrafo del 148. Sin consentimiento cambiaría el análisis al aborto forzado CPDF 146; el supuesto refiere solicitudes voluntarias \citep{cpdfS3Consolidado}.

\textit{El marco local y el criterio constitucional deben leerse conjuntamente.} CPDF 144--148 y CPEUM arts.~1 y 4 son el marco local y constitucional \citep{cpdfS3Consolidado,cpeumAulaS3}. La Primera Sala, en el AR 267/2023, consideró inconstitucionales porciones del régimen federal CPF 330--334 por afectar derechos reproductivos; no se debe convertir ese amparo en una afirmación no probada de derogación de todo el CPF ni ignorar el régimen local \citep{scjnAR2672023}. Ninguna de las dos solicitudes descritas permite inferir automáticamente responsabilidad penal: verificar los datos y garantizar atención respetuosa de derechos.
\end{casoestudio}

\subsection*{Síntesis de los encuadramientos}
\phantomsection\addcontentsline{toc}{subsection}{Síntesis de los encuadramientos}
El contraste separa los resultados de muerte de la lesión de Pedro y de las dos solicitudes aún no ejecutadas. El nexo causal es decisivo en los primeros, pero el parentesco conocido, las razones de género y la intervención en una decisión suicida agregan requisitos distintos. Una fractura se valora por su efecto en la salud; la interrupción del embarazo requiere precisar semanas, consentimiento y excluyentes. Ningún factor suple los restantes \citep{amuchateguiPenal2012,cpdfS3Consolidado,cpfOficial2026}.

La semejanza del verbo típico no borra las calidades específicas de sujetos y circunstancias; la divergencia entre códigos tampoco determina el fuero. Se requieren dictámenes para fractura y muerte, prueba del vínculo, antecedentes de violencia, comunicaciones y constancias clínicas. La perspectiva de género obliga a investigar sin convertir indicios en condena, y el control de las restricciones reproductivas exige distinguir petición de intervención realizada. Las dos resoluciones de la SCJN aportan criterios de investigación e interpretación, no sustitutos de esos datos \citep{scjnAR5542013,scjnAR2672023}.

\begin{landscape}
\thispagestyle{empty}
\phantomsection\addcontentsline{toc}{subsection}{Cuadro comparativo de los cinco casos}
\begingroup\scriptsize
\setlength{\tabcolsep}{4pt}
\renewcommand{\arraystretch}{1.62}
\begin{longtable}{@{}>{\raggedright\arraybackslash}p{0.12\linewidth}*{5}{>{\raggedright\arraybackslash}p{0.158\linewidth}}@{}}
\toprule
 & \textbf{1. Carlos / Roberto} & \textbf{2. Miguel / padre} & \textbf{3. Ana / expareja} & \textbf{4. Persona / conocido} & \textbf{5. Solicitudes en CDMX} \\
\midrule\endhead
Figura por examinar & Lesiones de Pedro; homicidio de la segunda víctima si se prueba el nexo. & Homicidio por parentesco conocido; contrastar con simple. & Violencia de pareja previa y posible feminicidio; contrastar con homicidio. & Inducción y auxilio al suicidio; no ejecución directa según relato. & 10 semanas: fuera de definición penal local; 16 semanas: excluyente por violación. \\
Conducta y resultado & Carlos / Pedro: golpe y fractura. Roberto / otra persona: agresión y muerte posterior. & Miguel / padre: agresión y muerte; filiación conocida. & Expareja / Ana: amenazas, control económico, agresiones y muerte. & Conocido / persona fallecida: persuasión y medios; actos finales propios. & Dos mujeres solicitan atención médica; no se presume procedimiento ya ejecutado. \\
Bien protegido & Integridad y salud de Pedro; vida de la segunda víctima. & Vida del ascendiente; el vínculo añade requisito típico. & Vida de Ana y protección ante violencia por razones de género. & Vida de quien recibe ayuda o inducción. & Autonomía reproductiva, salud y bienes tutelados según la norma aplicable. \\
Elemento diferenciador & Dictamen sobre fractura, nexo y dolo; discutir no es riña ni agravante. & Filiación y conocimiento del vínculo, nexo y dolo de matar. & Razones legales de género y contexto de agresiones previas. & Peso causal del auxilio o inducción frente a actos propios. & Plazo y consentimiento; violación excluye responsabilidad según norma local. \\
Contraste normativo & Ambos tutelan vida e integridad; el CPDF 130 gradúa lesiones por secuela o sanación y el CPF 289--293 distingue otros resultados. & Ambos exigen vínculo conocido; CPDF 125 y CPF 323 enumeran relaciones con alcances propios. & CPDF 148 Bis y CPF 325 exigen razones de género; la norma local detalla vínculos y contextos distintos. & CPDF 142 distingue inducción, ayuda y no consumación; CPF 312 reúne inducción y auxilio en el mismo supuesto inicial. & CPDF 144 fija el umbral de doce semanas; CPF 329 usa otra definición, a la luz del AR 267/2023. \\
Texto local / federal & CPDF 130: ``daño o alteración en su salud''; 123: ``prive de la vida a otro''. CPF 288: ``toda alteración en la salud''; 302: ``priva de la vida a otro''. & CPDF 125 y CPF 323 exigen ``con conocimiento de esa relación''; sin ese saber hay que revisar homicidio simple. & CPDF 148 Bis: ``por razones de género, prive de la vida a una mujer''; CPF 325: ``por una razón de género''. & CPDF 142: ``ayude a otro para que se prive de la vida''; CPF 312: ``prestare auxilio o indujere a otro para que se suicide''. & CPDF 144: ``después de la décima segunda semana de gestación''; CPF 329: ``en cualquier momento de la preñez''. \\
Prueba y alternativa & Pericia sobre fractura, necropsia, historia clínica y nexo; explorar causa independiente de muerte. & Acta de filiación, conocimiento probado y peritaje causal; descartar simple o lesión sin dolo de matar. & Mensajes, denuncias y pericias de contexto; investigar homicidio sin presuponer feminicidio. & Cronología de mensajes, medios aportados y capacidad; distinguir decisión previa y ejecución ajena. & Constancias de semanas, consentimiento y origen del embarazo; no equiparar solicitud con aborto realizado. \\
Error a evitar & Equiparar fractura con peligro vital o sucesión temporal con nexo causal probado. & Dar por demostrada la intención homicida solo por el parentesco. & Confundir deber de investigar con prueba automática de feminicidio o concurso. & Considerar toda persuasión inducción decisiva o ayuda como ejecución material. & Inferir aborto consumado a partir de la petición o aplicar el régimen federal sin competencia. \\
Fundamento / criterio judicial & CPDF 123--124, 130, 138; CPF 288, 302--305. & CPDF 125 (simple: 123); CPF 323. & CPDF 148 Bis, 200; CPF 325; SCJN AR 554/2013. & CPDF 142--143; CPF 312--313. & CPDF 144--148; CPEUM 1 y 4; SCJN AR 267/2023. \\
\bottomrule
\end{longtable}
\endgroup
\end{landscape}
\clearpage
\section{Conclusiones}
Considero que el encuadramiento debe tratarse como una hipótesis revisable: la vida, la integridad y la autonomía no quedan protegidas cuando se convierte un dato aislado (una fractura, una relación familiar, una amenaza o una semana de gestación) en una condena. El nexo causal separa las lesiones de la muerte atribuible; la calidad del sujeto y su conocimiento distinguen el parentesco; y la intervención de un tercero no equivale a que este haya ejecutado el suicidio. Confundir esas condiciones debilita la protección de las víctimas y el control del poder punitivo.

Sostengo que, ante la muerte de Ana, la ruta más responsable es investigar los antecedentes y las razones de género sin prejuzgar culpabilidad. En materia de embarazo, la semana de gestación importa, pero no desplaza el consentimiento ni la excluyente por violación: ignorarlos transforma un análisis de derechos en una conclusión penal errónea. Por eso daría prioridad a peritajes sobre fractura y muerte, documentos del vínculo familiar, registros de violencia y constancias clínicas; determinaría la ley temporal y el fuero antes de proponer sanciones.

La teoría del caso debe ser revisable cuando cambia la evidencia: la formulación de la imputación expone una calificación jurídica preliminar, mientras que la acusación posterior debe precisar hechos, clasificación y medios de prueba \citep[arts.~211, 311 y 335]{cnppOficial2025}. En ese tránsito, un dictamen puede desplazar lesiones hacia homicidio o descartarlo por falta de nexo; un antecedente corroborado puede modificar la investigación de la muerte de Ana. Mi posición es que corregir fundadamente el encuadramiento protege tanto el derecho de defensa como el esclarecimiento del hecho: aferrarse a la primera hipótesis permitiría que una inferencia ocupara el lugar de un elemento típico.\footnote{Se utilizó GitHub Copilot para organizar conceptos, contrastar fuentes y elaborar un borrador de redacción y maquetación. Corresponde al estudiante revisar, corregir y asumir el análisis antes del envío.}

\clearpage
\begingroup\footnotesize
\begin{thebibliography}{7}
\bibitem[Amuchategui Requena(2012)]{amuchateguiPenal2012}
Amuchategui Requena, I.~G. (2012). \emph{Derecho penal} (4.ª ed.). Oxford University Press México.

\bibitem[Cámara de Diputados(2025)]{cnppOficial2025}
Cámara de Diputados del H. Congreso de la Unión. (2025). \emph{Código Nacional de Procedimientos Penales} (última reforma indicada: DOF 28-11-2025). \url{https://www.diputados.gob.mx/LeyesBiblio/pdf/CNPP.pdf}

\bibitem[Cámara de Diputados(2026a)]{cpfOficial2026}
Cámara de Diputados del H. Congreso de la Unión. (2026a). \emph{Código Penal Federal} (última reforma indicada: DOF 13-03-2026). \url{https://www.diputados.gob.mx/LeyesBiblio/pdf/CPF.pdf}

\bibitem[Cámara de Diputados(2026b)]{cpeumAulaS3}
Cámara de Diputados del H. Congreso de la Unión. (1917/2026b). \emph{Constitución Política de los Estados Unidos Mexicanos} (última reforma indicada: DOF 07-10-2026). \url{https://www.diputados.gob.mx/LeyesBiblio/pdf/CPEUM.pdf}

\bibitem[Consejería Jurídica de la CDMX(2026)]{cpdfS3Consolidado}
Consejería Jurídica y de Servicios Legales de la Ciudad de México. (2026). \emph{Código Penal para el Distrito Federal} (texto consolidado; última reforma indicada: G.O. CDMX 21-08-2026). \url{https://data.consejeria.cdmx.gob.mx/images/leyes/2025/2026/50226/CODIGO_PENAL_PARA_EL_DF_12.9.pdf}

\bibitem[Suprema Corte de Justicia de la Nación(2015)]{scjnAR5542013}
Suprema Corte de Justicia de la Nación. (2015, 25 de marzo). \emph{Amparo en revisión 554/2013} (Primera Sala). \url{https://www.scjn.gob.mx/derechos-humanos/sites/default/files/sentencias-emblematicas/sentencia/2020-12/AR%20554-2013.pdf}

\bibitem[Suprema Corte de Justicia de la Nación(2023)]{scjnAR2672023}
Suprema Corte de Justicia de la Nación. (2023). \emph{Amparo en revisión 267/2023} (Primera Sala). \url{https://www2.scjn.gob.mx/juridica/engroses/1/2023/2/2_311450_6648_firmado.pdf}
\end{thebibliography}
\endgroup
\end{document}